\begin{table*}[t]
\centering
\caption{How far each method's mean ARI moves with the candidate design and the scorer. Means over the nine small datasets of the per-dataset medians, confirmatory seeds 3 to 9. The lowest, middle and highest setting per count and every setting differ only for Leiden, since $k$ has one setting per count. S1 scores the same partitions on all 100 subsamples. S3 is chooseR's own R code on all 100, with a point interval where its bootstrap gives none (the protocol's primary repair). chooseR has only the finest-stable reading, and the most-stable reading uses no interval, so it equals S1.}
\label{tab:designs}
\footnotesize
\setlength{\tabcolsep}{4pt}
\begin{tabular}{lcccccc}
\toprule
 & \multicolumn{3}{c}{Finest stable} & \multicolumn{3}{c}{Most stable} \\
\cmidrule(lr){2-4}\cmidrule(lr){5-7}
Design and scorer & MICA & Leiden & $k$-means & MICA & Leiden & $k$-means \\
\midrule
Original grids & 0.647 & 0.662 & 0.452 & 0.452 & 0.650 & 0.374 \\
Matched, middle setting & 0.647 & 0.591 & 0.448 & 0.452 & 0.454 & 0.370 \\
Matched, lowest setting & 0.647 & 0.563 & 0.448 & 0.452 & 0.445 & 0.370 \\
Matched, highest setting & 0.647 & 0.579 & 0.448 & 0.452 & 0.481 & 0.370 \\
Matched, every setting & 0.647 & 0.569 & 0.448 & 0.452 & 0.420 & 0.370 \\
Matched, 100 subsamples (S1) & 0.679 & 0.528 & 0.434 & 0.486 & 0.428 & 0.404 \\
Original, chooseR's code (S3) & 0.676 & 0.608 & 0.457 & n/a & n/a & n/a \\
Matched, chooseR's code (S3) & 0.679 & 0.591 & 0.446 & n/a & n/a & n/a \\
\bottomrule
\end{tabular}
\end{table*}
